\section{Discovery}\label{sec:construction}

Chen's interpretation of noncongruence modular curves as Hurwitz spaces
was the starting point~\cite[\S5.4]{Chen2025}. On his distinguished
component, one positive-rank rational fiber would suffice
\cite[Proposition~5.6 and the following remark]{Chen2025}, but our
computation found none. We instead searched for three-point
$\Aut(\Sz(8))$-covers whose cyclic cubic quotient could yield $\Sz(8)$
after base change; an order-three symmetry of a $(5,5,5)$ cover suggested
this route.

\subsection{The three passports with outer inertia of order three}\label{sec:passports}

For a three-point regular $G$-cover, the two outer branch points form one
Galois orbit. The proof of \cref{prop:cyclic-prime} shows that their common
residue field is exactly $\Q(\zeta_3)$. The third branch point is rational.
A fractional linear transformation over $\Q$ places the third point at
$\infty$ and the conjugate pair at $\zeta_3,\zeta_3^2$.

We restrict the two outer inertia orders to $3$. There is one order-$3$
class $D$ in each outer coset, and the two classes are inverses. The class
in the rational branch slot must be rational in $G$. The order-$7$
classes of $N$ fuse in $G$, as do the order-$13$ classes; together with the
order-$5$ class they give the following possibilities. The two order-$4$
classes are inverse and remain distinct, so neither is rational. The
involution class, of type $1\,2^{32}$, would give the impossible genus
$1-65+(32+40+40)/2=-8$. These observations make the table exhaustive
under the stated order-$3$ restriction.

\begin{center}
\begin{tabular}{@{}cccc@{}}
\toprule
Order in the rational slot&Cycle type&Source genus&Nielsen classes\\
\midrule
$5$&$5^{13}$&$2$&$3$\\
$7$&$1^2 7^9$&$3$&$6$\\
$13$&$13^5$&$6$&$4$\\
\bottomrule
\end{tabular}
\end{center}
The other two cycles have type $1^5 3^{20}$. For example, the last row has
genus $1-65+(60+40+40)/2=6$.

The counts are for generating product-one triples in $(C,D,D^{-1})$
with ordered branch slots, modulo simultaneous conjugation by $G$.
Fixing the first entry $c\in C$ leaves respectively $45$, $42$, and $52$
choices of the second entry that generate $G$ and give the required third
entry. The centralizers of $c$ have orders $15$, $7$, and $13$.
Their action is free because a generating triple has centralizer $Z(G)=1$;
division gives $3$, $6$, and $4$. Thus the $(13,3,3)$ passport is not rigid.
Complex conjugation fixes two of its four classes and interchanges the
other two. Reality alone does not determine which classes descend to $\Q$;
the polynomial certificate proves descent for the class used here.

The order-$5$ construction supplied a genus-two cover over
$\Q(\theta)$ with $\theta^3-\theta^2+2\theta+2=0$. The order-$13$ row
supplied the rational cover.

\subsection{Numerical construction and arithmetic coordinates}

The triangle-group method associates a finite-index subgroup of the
triangle group of signature $(13,3,3)$ to a permutation triple. Expansions
of holomorphic differentials give a canonical model of the source curve.
For the chosen cover the canonical model is a genus-six curve in $\PP^5$,
realized as a quadric section of a smooth del Pezzo surface of degree $5$.
In particular, it is neither hyperelliptic, trigonal, nor a plane quintic.
The canonical hyperplane divisor $H$ has degree $10$.

A marking of the del Pezzo surface is additional arithmetic data.  The
descent computation identified the Galois action on its five degree-four
pencils with the full group $S_5$.  Coordinates adapted to an ordered
marking therefore need not be rational even when the curve and covering
map descend to $\Q$.  Invariant data and the branch fibers produced an
exact rational model, after which the map was represented by a ratio of
sections of $\OO_Z(8H)$.  Some entries of the exact kernel vectors defining
those sections have $480038$ decimal digits.  This size motivated a
different generator of the function field rather than further manipulation
of the same coordinates.

\subsection{Generators with poles over infinity}

The pole fiber is $13(P_0+P_4)$, with the rational point $P_0$ and
degree-four point $P_4$ identified in \cref{sec:infinity}.
In the discovery model the first non-gap at $P_0$
is $7$, so $P_0$ is not a Weierstrass point. Choose a nonconstant function
$x\in L(7P_0)$. The coefficient bound is \cref{lem:polebound}.

For $x\in L(7P_0)$ the list consists of thirteen entries $7/13$ and
fifty-two zeros, recovering
$\deg_t c_i\leq\lfloor7\min(i,13)/13\rfloor\leq7$.
The function $x$ generates $\Q(Z)/\Q(t)$ because the permutation action
is primitive: a proper intermediate field would give a block system,
while a function with only the pole $P_0$ cannot belong to $\Q(t)$.
Trace-centering kills the coefficient of $X^{64}$.
The resulting $f$ has primitive integer height $249$ digits. Its rational
coefficients have numerators of at most $244$ digits and common denominator
$13^{182}$; these are separate coefficient statistics.

The same search in $L(3P_0+P_4)$ produced $F_3$. The pole bound and
exact bridge are given in \cref{sec:smaller}.

\subsection{Recovering the linear series without eliminating the base locus}

Write $t=N/D$ with $N,D\in V=H^0(Z,\OO_Z(8H))$, so $\dim V=75$.
Their common base divisor $R$ has degree $15$, and
$\operatorname{div}(D)=13(P_0+P_4)+R$. The condition that $A/D$ belong to
$L(7P_0)$ consists of vanishing of $A$ along $R$, to order $13$ along
$P_4$, and to order $6$ at $P_0$.

The base-divisor condition can be imposed by multiplication of sections.
Indeed,
\begin{equation}\label{eq:base-locus-linear}
 NV+DV=H^0(Z,\OO_Z(16H)(-R)).
\end{equation}
The right-hand side has dimension $145+1-6=140$. To compute the dimension
on the left, remove the common divisor from $N$ and $D$. The resulting
pencil is base-point free. Its relations identify $NV\cap DV$ with
$H^0(Z,\OO_Z(R))$, which has dimension $15+1-6=10$. Hence the left-hand
side also has dimension $150-10=140$, proving \eqref{eq:base-locus-linear}.
Since $V$ is globally generated,
\begin{equation}\label{eq:membership}
 A\text{ vanishes along }R\quad\Longleftrightarrow\quad AV\subseteq NV+DV.
\end{equation}
For the reverse implication, choose locally a section in $V$ that is a
unit along each point of $R$ and apply \eqref{eq:base-locus-linear}.

This replaces explicit equations for the base locus by a linear-algebra
test. Sections of $\OO_Z(16H)$ are determined by jets of length $34$ at
the five geometric pole points: a section vanishing on all those jets
would have at least $170$ zeros, exceeding its degree $160$.
The same jets therefore suffice to implement \eqref{eq:membership}.

The Riemann--Roch calculation was performed at $92$ primes used for Chinese
remaindering and at eight held-out primes.  At every prime the computed
space $L(7P_0)$ had dimension two and pole-order flag $(7,0)$; all $473$
coefficient slots were reconstructed, and an exact rational control
identified the result with an affine change of $f$.  Thus the calculation
also supplies the stated first non-gap.  The same run for
$L(3P_0+P_4)$ gave dimension two, pole-order flag $(3,0)$, and all $319$
slots reconstructed with the same held-out checks.  The resultant
calculation in \cref{app:smaller} then proves over $\Q(t)$ that its reduced
polynomial has the same function field as $f$.

Choosing these generators turns the large section representation into a
compact plane equation. The resulting singularities account for the
apparent branch points treated in \cref{sec:certificate}.
